% Section 1.1: Evidence for Dark Matter
% Structure: Mini-Midi-Maxi classification (Cirelli et al. 2024)

\section{Evidence for Dark Matter}
\label{sec:1.1}
\aure{stopped here: working in review mode}

The case for dark matter rests not on a single observation but on a convergence of independent lines of evidence spanning a wide range of astrophysical scales.
From the kinematics of individual galaxies, through the dynamics and gravitational lensing of galaxy clusters, to the precision cosmology of the cosmic microwave background, each class of observation independently requires a dominant non-luminous matter component.
We follow the hierarchical classification of Cirelli et al.~\cite{Cirelli:2024ssz}, organizing the evidence from galactic (``mini'') to cluster (``midi'') to cosmological (``maxi'') scales.
The internal consistency of these measurements --- obtained with different techniques, probing different environments --- is probably the best argument that dark matter is a physical reality rather than an artifact of any single observational systematic.

\subsection{Galactic Scale}
\label{sec:1.1.1}

The most direct evidence for dark matter comes from the internal kinematics of galaxies.
In a spiral galaxy, the circular velocity $v_c(r)$ of stars and gas at distance $r$ from the center is determined by the enclosed mass through $v_c(r) = \sqrt{GM(r)/r}$.
If the visible mass --- concentrated in a central bulge and an exponential disk --- were the only source of gravity, the rotation curve should decline as $v_c \propto r^{-1/2}$ beyond the luminous extent, following the standard Keplerian falloff~\cite{Cirelli:2024ssz}.

Observations tell a different story.
The first hints of anomalous galactic dynamics appeared with Babcock's 1939 study of M31, though the measurements lacked the radial extent and precision to be conclusive~\cite{Cirelli:2024ssz}.
The modern picture took shape in 1970, when Rubin and Ford~\cite{Rubin:1970zza} obtained the first high-precision optical rotation curve of M31, tracing approximately 70 H\,\textsc{ii} regions out to ${\sim}\,22$~kpc and finding no evidence of the expected Keplerian decline.
In the same year, Freeman identified similar discrepancies in the disk dynamics of several spiral galaxies~\cite{Cirelli:2024ssz}.
Over the following decades, radio observations of the 21-cm hydrogen line extended rotation curves well beyond the optical disk, establishing that $v_c(r)$ remains approximately constant --- or even rises slightly --- out to the largest measurable radii in hundreds of systems.
To maintain a constant circular velocity, the enclosed mass must grow linearly with radius, $M(r) \propto r$, which requires an extended, roughly spherical halo of invisible matter with density $\rho(r) \propto r^{-2}$ at large radii~\cite{Cirelli:2024ssz}.
The comprehensive analysis of Persic, Salucci, and Stel~\cite{Persic:1995ru} synthesized rotation curves across a wide range of galaxy luminosities into a ``Universal Rotation Curve'', demonstrating that the dark-to-luminous mass ratio increases systematically with decreasing galaxy luminosity.

Modern data have only strengthened this picture.
The Spitzer Photometry and Accurate Rotation Curves (SPARC) compilation~\cite{Lelli:2016zqa} provides high-quality mass models and rotation curves for 175 disk galaxies, spanning a broad range of surface brightnesses and morphologies.
SPARC has confirmed the ubiquity of flat rotation curves and revealed tight empirical scaling relations --- notably the baryonic Tully-Fisher relation, $v_\infty \propto M_b^{1/4}$, linking the asymptotic rotation velocity to the total baryonic mass --- that any successful dark matter model must reproduce~\cite{Cirelli:2024ssz}.
More recently, high-precision astrometric data from Gaia Data Release~3 have been used to map the Milky Way's rotation curve out to ${\sim}\,100$~kpc.
Analyses by Jiao et al.~\cite{Jiao:2023aci} and Ou et al.~\cite{Ou:2023adg} report the first hints of a Keplerian decline in the Galaxy's far outskirts, placing new constraints on the total mass and spatial extent of the Milky Way's dark matter halo~\cite{Cirelli:2024ssz}.
Figure~\ref{fig:rotation_curves} illustrates this progression, from the original optical rotation curve of M31 through modern compilations spanning hundreds of galaxies to individual mass decompositions that reveal the dominant dark matter contribution at large radii.

%
% \includegraphics[width=0.45\textwidth,height=0.22\textheight]{figures/RotationCurveMW.pdf} \\[2mm]
% \includegraphics[width=0.45\textwidth,height=0.23\textheight]{figures/rotationcurve3_compilation_fromSofue1999.pdf}$\qquad$%
% \includegraphics[width=0.44\textwidth]{figures/RotationCurveCompilationSPARC.pdf}\\[2mm]
% \includegraphics[width=0.97\textwidth]{figures/RotationCurveSample.pdf}
\begin{figure}[t]
	\centering
	\includegraphics[width=0.45\textwidth,height=0.22\textheight]{figures/rubin-m31.pdf}
	\caption{\aure{update caption} Rotation curves of spiral galaxies. \textit{Top left}: the original rotation curve of M31 by Rubin \& Ford (1970). \textit{Top right}: a recent rotation curve for the Milky Way from Jiao et al.~\cite{Jiao:2023aci}, showing a hint of Keplerian decline. \textit{Center left}: a compilation of ${\sim}\,50$ galaxies from Sofue et al.~(1999). \textit{Center right}: the SPARC compilation of 175 galaxies~\cite{Lelli:2016zqa}, color-coded by surface brightness (red is denser, blue is fainter). \textit{Bottom left}: a denser galaxy, dominated in its center by baryons, mostly in stars. \textit{Bottom right}: a fainter galaxy, dominated by a dark component, with most baryons as gas.
		Both panels show the total rotation velocity (black), disk stars (dotted magenta), bulge (dot-dashed red), gas (dashed blue), and inferred dark matter halo (grey).
		Figure from Cirelli et al.~\cite{Cirelli:2024ssz}.
	}
	\label{fig:rotation_curves}
\end{figure}

An independent line of evidence comes from dispersion-supported systems.
Dwarf spheroidal galaxies (dSphs) --- small, gas-poor satellites of the Milky Way such as Draco, Sculptor, and Segue~1 --- lack ordered rotation, so their masses are instead inferred from the line-of-sight velocity dispersions of their constituent stars.
For a dispersion-supported system, the mass enclosed within the three-dimensional half-light radius is well approximated by $M(<r_{1/2}) \approx 3\sigma_*^2 r_{1/2} / G$~\cite{Wolf:2009tu}.
Applying this estimator reveals that dSphs are extremely dark-matter dominated, with mass-to-light ratios ranging from $M/L \sim 10\,M_\odot/L_\odot$ in the classical dwarfs to several hundred in Draco ($M/L \sim 300$) and well over a thousand in ultra-faint systems like Segue~1~\cite{Walker:2009zp,Cirelli:2024ssz,Bullock:2017xww}.
These extreme ratios are difficult to explain through any baryonic mechanism and rank dSphs among the most dark-matter-dominated objects known.
Their low astrophysical backgrounds also make them prime targets for indirect dark matter searches, a point to which we return in Section~\ref{sec:1.3}.

\subsection{Galaxy Cluster Scale}
\label{sec:1.1.2}

Moving to larger scales, the dynamics of galaxy clusters provide a second, independent class of evidence for dark matter.
The earliest quantitative argument dates to 1933, when Zwicky~\cite{Cirelli:2024ssz} studied the radial velocities of galaxies in the Coma Cluster and applied the virial theorem --- which relates the time-averaged kinetic energy of a gravitationally bound system to its potential energy, $2\langle K \rangle + \langle U \rangle = 0$ --- to estimate the cluster's total mass.
He found a mass-to-light ratio of order ${\sim}\,400$, far exceeding what could be accounted for by the luminous galaxies alone, and concluded that the cluster must contain a large amount of unseen ``dunkle Materie''~\cite{Cirelli:2024ssz}.

Modern X-ray observations have placed these early dynamical estimates on much firmer ground.
Galaxy clusters contain massive reservoirs of ionized gas --- the intracluster medium (ICM) --- heated to temperatures of $10^7$--$10^8$~K as it falls into the cluster potential well.
This gas radiates thermally through bremsstrahlung emission, observed in the X-ray band~\cite{Cirelli:2024ssz}.
Under the assumption of hydrostatic equilibrium, the outward gas pressure gradient must balance the inward gravitational pull, allowing the total gravitating mass to be reconstructed from the measured gas temperature and density profiles.
These analyses consistently show that the visible galaxies plus the X-ray-emitting gas account for only ${\sim}\,15\%$ of the total cluster mass~\cite{Cirelli:2024ssz}, with the rest attributable to a dark, non-luminous component.

The Sunyaev-Zel'dovich (SZ) effect offers a complementary probe.
When cosmic microwave background (CMB) photons traverse a cluster, they undergo inverse-Compton scattering off the hot ICM electrons, gaining energy and producing a characteristic spectral distortion~\cite{Cirelli:2024ssz,Bullock:2017xww}.
The amplitude of this distortion is proportional to the integrated electron pressure along the line of sight and, unlike X-ray surface brightness, does not diminish with redshift.
Combined with independent constraints on the baryon density from Big Bang Nucleosynthesis (BBN) and the CMB, SZ measurements of the cluster baryon fraction $f_b \approx \Omega_b / \Omega_m$ yield a total matter density $\Omega_m \approx 0.3$, well above the baryonic contribution~\cite{Cirelli:2024ssz}.

Gravitational lensing provides a direct, model-independent method for mapping the total mass distribution of clusters, independent of assumptions about dynamical equilibrium.
Strong lensing in the dense cluster core produces highly distorted arcs and multiple images of background galaxies, while weak lensing measures the subtle statistical shearing of background galaxy shapes over larger angular scales~\cite{Cirelli:2024ssz,Bullock:2017xww}.
Both regimes allow the total mass distribution --- visible and dark --- to be reconstructed without reliance on assumptions about the dynamical state of the system.

The most striking application of gravitational lensing is the Bullet Cluster (1E\,0657-558), analyzed by Clowe et al.~\cite{Clowe:2006eq}.
This system consists of two galaxy clusters that passed through each other roughly 150 million years ago.
X-ray observations reveal that the collisional baryonic gas from the two clusters interacted strongly, producing a bow shock and lagging behind the galaxies.
Weak-lensing mass maps, however, show that the bulk of the total mass is spatially offset from the gas, coinciding instead with the collisionless galaxies~\cite{Clowe:2006eq}.
This spatial separation is clearly visible in Figure~\ref{fig:bullet_cluster}: the green mass contours reconstructed through weak lensing are clearly displaced from the X-ray-emitting gas, providing a direct, model-independent demonstration that the dominant mass component is collisionless --- exactly the behavior expected of dark matter and precisely the opposite of what modified gravity theories such as MOND would predict, since in those frameworks the lensing signal should trace the baryons~\cite{Cirelli:2024ssz}.
Harvey et al.~\cite{Harvey:2015hha} extended this analysis to 72 colliding clusters, confirming the separation of dark matter from baryonic gas at greater than $7\sigma$ significance and setting stringent upper limits on the dark matter self-interaction cross-section~\cite{Cirelli:2024ssz}.

\begin{figure}[t]
	\centering
	\includegraphics[width=0.6\textwidth]{figures/f1b_new_jpeg2ps.pdf}
	\caption{The Bullet Cluster (1E\,0657-558).
		The colored map shows the X-ray image of the hot baryonic gas from \textit{Chandra} observations, while the green contours trace the total mass distribution reconstructed through weak gravitational lensing.
		The clear spatial offset between the two demonstrates that the dominant mass component is collisionless, as expected for dark matter.
		The white bar corresponds to 200~kpc.
		Adapted from Clowe et al.~\cite{Clowe:2006eq}, figure from Cirelli et al.~\cite{Cirelli:2024ssz}.
	}
	\label{fig:bullet_cluster}
\end{figure}

\subsection{Cosmological Scale}
\label{sec:1.1.3}

At the largest scales, precision cosmological observations provide the most quantitatively precise evidence for dark matter and its non-baryonic nature.

The cosmic microwave background radiation --- the thermal afterglow of the hot, dense early Universe --- encodes detailed information about the matter content of the cosmos.
The temperature anisotropy power spectrum, measured with high precision by the Planck satellite, encodes the physics of acoustic oscillations in the primordial baryon-photon plasma before recombination~\cite{Aghanim:2018eyx,Cirelli:2024ssz}.
Baryons load the oscillating fluid, enhancing the compressional (odd-numbered) acoustic peaks relative to the rarefaction (even-numbered) ones; the ratio of the first to the second peak therefore directly constrains the physical baryon density~\cite{Cirelli:2024ssz}.
Dark matter, on the other hand, does not couple to the photon pressure and instead provides the gravitational scaffolding within which the oscillations occur.
Its abundance sets the overall depth of the potential wells and modulates the relative heights of successive peaks --- in particular, a universe without cold dark matter would produce a strongly suppressed third peak, in clear contradiction with the data~\cite{Cirelli:2024ssz}.
The Planck 2018 analysis yields $\Omega_c h^2 = 0.1200 \pm 0.0012$ for the cold dark matter density and $\Omega_b h^2 = 0.02237 \pm 0.00015$ for the baryon density~\cite{Aghanim:2018eyx}, confirming that dark matter is roughly five times more abundant than ordinary matter.
As CMB photons propagate toward us from the surface of last scattering, their paths are further deflected by intervening large-scale structure, an effect known as CMB gravitational lensing.
This lensing smooths the acoustic peaks at high multipoles and converts a fraction of the primordial E-mode polarization into B-modes~\cite{Cirelli:2024ssz,Planck:2018lbu}, providing an additional handle on the integrated matter distribution along the line of sight.
The dramatic impact of cold dark matter on both the CMB and matter power spectra is illustrated in Figure~\ref{fig:cmb_lss}: the bottom row compares the observed spectra with those predicted in a universe without dark matter, showing just how radically different the cosmic picture would be.

\begin{figure}[t]
	\centering
	\includegraphics[width=0.47\textwidth]{figures/CMB_PlanckMap.jpg}~%
	\includegraphics[width=0.47\textwidth,height=165pt]{figures/LSS_sdss_release14.jpg}\\[2mm]
	\includegraphics[width=0.47\textwidth]{figures/CMB_Planck2014_TT.pdf}~%
	\includegraphics[width=0.47\textwidth,height=170pt]{figures/LSS_powerspectrum.pdf}\\[2mm]
	\includegraphics[width=0.47\textwidth,height=170pt]{figures/CMBnoDMAS}~\qquad%
	\includegraphics[width=0.47\textwidth]{figures/LSSpowerDMAS}
	\caption{The imprint of dark matter on cosmological observables. \textit{Top left}: the Planck CMB temperature anisotropy map. \textit{Top right}: the large-scale distribution of galaxies from the SDSS survey. \textit{Middle row}: the CMB acoustic power spectrum (left) and the matter power spectrum (right) as observed. \textit{Bottom row}: the same quantities computed in a universe without cold dark matter, showing the dramatic suppression of the third acoustic peak and the absence of small-scale structure.
		Figure from Cirelli et al.~\cite{Cirelli:2024ssz}.
	}
	\label{fig:cmb_lss}
\end{figure}

The observed large-scale structure of the Universe delivers a complementary argument for dark matter.
In a purely baryonic universe, structure formation would have been severely hindered: before recombination, baryons were tightly coupled to photons through Thomson scattering, and their density perturbations were erased on small scales by photon diffusion --- a process known as Silk damping~\cite{Cirelli:2024ssz}.
Cold dark matter, which does not interact electromagnetically, was immune to this damping and began to cluster gravitationally during the matter-dominated era, well before recombination.
The potential wells established by dark matter then served as the seeds into which baryons fell after decoupling, seeding the formation of galaxies and clusters~\cite{Cirelli:2024ssz}.
Without this pre-existing dark matter framework, the observed matter power spectrum $P(k)$ cannot be reproduced.

Baryon acoustic oscillations (BAO) preserve the imprint of the primordial sound waves frozen into the matter distribution at recombination.
The characteristic comoving scale of ${\sim}\,150$~Mpc, set by the sound horizon at decoupling, appears as a well-defined peak in the two-point correlation function of galaxies and provides a ``standard ruler'' for measuring the expansion history~\cite{Cirelli:2024ssz}.
Its first detection in the galaxy distribution by Eisenstein et al.~\cite{SDSS:2005xqv} delivered independent confirmation of the $\Lambda$CDM parameters.
Observations of Type~Ia supernovae as standardizable candles, meanwhile, established the accelerated expansion of the Universe~\cite{SupernovaSearchTeam:1998fmf,SupernovaCosmologyProject:1998vns}.
Combined with BAO and CMB data, these measurements constrain the dark energy density to $\Omega_\Lambda \approx 0.7$ and the total matter density to $\Omega_m \approx 0.3$~\cite{Cirelli:2024ssz,Pan-STARRS1:2017jku}.

The final, logically essential element is Big Bang Nucleosynthesis.
The production of light elements during the first few minutes after the Big Bang is exquisitely sensitive to the baryon-to-photon ratio.
Measured primordial abundances of deuterium, $^4$He, and $^7$Li tightly constrain the baryon density to $\Omega_b h^2 \approx 0.022$~\cite{Cirelli:2024ssz}, in excellent agreement with the independent CMB determination.
Since the total matter density is $\Omega_m h^2 \approx 0.14$ --- nearly seven times larger --- the majority of the matter in the Universe cannot be made of baryons~\cite{Cirelli:2024ssz}.
This is the logical bridge: the astrophysical evidence reviewed in this section demonstrates that dark matter exists; BBN proves that it is not composed of ordinary atoms.
Whatever dark matter is, it must be a new form of matter beyond the Standard Model --- a conclusion that motivates the particle physics discussion of the following section.
